\documentclass[10pt,a4paper]{amsart}
\usepackage{amsmath,amssymb,amsthm,graphicx}
% Retained minimal preamble for the Stacks Project chapter more-morphisms.tex (More on Morphisms).
% Verbatim excerpt of the chapter preamble: Xy-pic setup (native preamble lines 14-18) and the native
% theorem declarations with the mathematical macros (native preamble lines 154-217). Copied mechanically;
% nothing added, removed or redefined. The wrapper class and the amsmath/amssymb/amsthm/graphicx base
% packages are supplied by the assembler, so the class/externaldocument part of the native preamble is omitted.
% For commutative diagrams we use Xy-pic
\usepackage[all]{xy}

% We use 2cell for 2-commutative diagrams.
\xyoption{2cell}
% Theorem environments.
%
\theoremstyle{plain}
\newtheorem{theorem}[subsection]{Theorem}
\newtheorem{proposition}[subsection]{Proposition}
\newtheorem{lemma}[subsection]{Lemma}

\theoremstyle{definition}
\newtheorem{definition}[subsection]{Definition}
\newtheorem{example}[subsection]{Example}
\newtheorem{exercise}[subsection]{Exercise}
\newtheorem{situation}[subsection]{Situation}

\theoremstyle{remark}
\newtheorem{remark}[subsection]{Remark}
\newtheorem{remarks}[subsection]{Remarks}

\numberwithin{equation}{subsection}

% Macros
%
\def\lim{\mathop{\mathrm{lim}}\nolimits}
\def\colim{\mathop{\mathrm{colim}}\nolimits}
\def\Spec{\mathop{\mathrm{Spec}}}
\def\Hom{\mathop{\mathrm{Hom}}\nolimits}
\def\Ext{\mathop{\mathrm{Ext}}\nolimits}
\def\SheafHom{\mathop{\mathcal{H}\!\mathit{om}}\nolimits}
\def\SheafExt{\mathop{\mathcal{E}\!\mathit{xt}}\nolimits}
\def\Sch{\mathit{Sch}}
\def\Mor{\mathop{\mathrm{Mor}}\nolimits}
\def\Ob{\mathop{\mathrm{Ob}}\nolimits}
\def\Sh{\mathop{\mathit{Sh}}\nolimits}
\def\NL{\mathop{N\!L}\nolimits}
\def\CH{\mathop{\mathrm{CH}}\nolimits}
\def\proetale{{pro\text{-}\acute{e}tale}}
\def\etale{{\acute{e}tale}}
\def\QCoh{\mathit{QCoh}}
\def\Ker{\mathop{\mathrm{Ker}}}
\def\Im{\mathop{\mathrm{Im}}}
\def\Coker{\mathop{\mathrm{Coker}}}
\def\Coim{\mathop{\mathrm{Coim}}}

% Boxtimes
%
\DeclareMathSymbol{\boxtimes}{\mathbin}{AMSa}{"02}

%
% Macros for moduli stacks/spaces
%
\def\QCohstack{\mathcal{QC}\!\mathit{oh}}
\def\Cohstack{\mathcal{C}\!\mathit{oh}}
\def\Spacesstack{\mathcal{S}\!\mathit{paces}}
\def\Quotfunctor{\mathrm{Quot}}
\def\Hilbfunctor{\mathrm{Hilb}}
\def\Curvesstack{\mathcal{C}\!\mathit{urves}}
\def\Polarizedstack{\mathcal{P}\!\mathit{olarized}}
\def\Complexesstack{\mathcal{C}\!\mathit{omplexes}}
% \Pic is the operator that assigns to X its picard group, usage \Pic(X)
% \Picardstack_{X/B} denotes the Picard stack of X over B
% \Picardfunctor_{X/B} denotes the Picard functor of X over B
\def\Pic{\mathop{\mathrm{Pic}}\nolimits}
\def\Picardstack{\mathcal{P}\!\mathit{ic}}
\def\Picardfunctor{\mathrm{Pic}}
\def\Deformationcategory{\mathcal{D}\!\mathit{ef}}
\newenvironment{reference}{}{}
\newenvironment{slogan}{}{}
\newenvironment{history}{}{}
\usepackage{cleveref}
\crefname{originallemma}{original lemma}{original lemmas}
\Crefname{originallemma}{Original lemma}{Original lemmas}
\makeatletter
\newlabel{more-algebra-definition-regular}{{more-algebra-definition-regular}{source}{}{sourcepointer-1}{}}
\newlabel{more-algebra-definition-regular@cref}{{[originallemma][0][]more-algebra-definition-regular}{[1][0][]source}{}{}{}}
\makeatother
\crefname{originallemma}{original lemma}{original lemmas}
\Crefname{originallemma}{Original lemma}{Original lemmas}
\makeatletter
\newlabel{morphisms-lemma-flat-characterize}{{morphisms-lemma-flat-characterize}{source}{}{sourcepointer-2}{}}
\newlabel{morphisms-lemma-flat-characterize@cref}{{[originallemma][0][]morphisms-lemma-flat-characterize}{[1][0][]source}{}{}{}}
\makeatother
\crefname{originallemma}{original lemma}{original lemmas}
\Crefname{originallemma}{Original lemma}{Original lemmas}
\makeatletter
\newlabel{varieties-lemma-geometrically-regular}{{varieties-lemma-geometrically-regular}{source}{}{sourcepointer-3}{}}
\newlabel{varieties-lemma-geometrically-regular@cref}{{[originallemma][0][]varieties-lemma-geometrically-regular}{[1][0][]source}{}{}{}}
\makeatother
\title{Five equivalent formulations of regularity for a morphism of schemes}
\author{The Stacks Project Authors}
\date{stacks-project:more-morphisms.tex:local-snapshot-20261011; GNU FDL 1.2 or later}
\begin{document}
\maketitle
\noindent\textit{Editorial note.} Statement span is lines 5681-5699 of sources/stacks-more-morphisms/pinned-original-source.tex; proof span is lines 5701-5733, the complete author proof (31 body lines).\par\medskip
\noindent\textit{Editorial note.} Navigation stubs: 3 cited-but-not-included labels are supplied as non-load-bearing pointers with their pinned tags where the snapshot tag index has them.\par\medskip
\noindent\textit{Editorial note.} extra\_packages is empty on purpose: the retained preamble already loads xy with options and a second declaration clashes in the assembler wrapper.\par\medskip
\noindent\textit{Editorial note.} Text is retained verbatim, including the author's own typos and placeholders.\par\medskip
\noindent\textit{Allowed standard prerequisites (editorial).}\begin{itemize}
\item Explicit prerequisite boundary. Regular morphisms, geometrically regular fibres, regular ring maps, and affine open covers; the reduction is retained.
\end{itemize}\medskip
\ifdefined\hypertarget\hypertarget{sourcepointer-1}{}\fi
\noindent\textit{Source reference.} Original lemma more-algebra-definition-regular: sources/stacks-more-morphisms/pinned-original-source.tex, source lines 5733--5733. Cited by the retained statement or proof but not part of the collected spans; supplied only as a navigation pointer. The argument does not depend on its content.\par\medskip
\ifdefined\hypertarget\hypertarget{sourcepointer-2}{}\fi
\noindent\textit{Source reference.} Original lemma morphisms-lemma-flat-characterize: sources/stacks-more-morphisms/pinned-original-source.tex, source lines 5733--5733. Cited by the retained statement or proof but not part of the collected spans; supplied only as a navigation pointer. The argument does not depend on its content.\par\medskip
\ifdefined\hypertarget\hypertarget{sourcepointer-3}{}\fi
\noindent\textit{Source reference.} Original lemma varieties-lemma-geometrically-regular: sources/stacks-more-morphisms/pinned-original-source.tex, source lines 5733--5733. Cited by the retained statement or proof but not part of the collected spans; supplied only as a navigation pointer. The argument does not depend on its content.\par\medskip
\begin{lemma}
\label{lemma-regular}
Let $f : X \to Y$ be a morphism of schemes.
Assume all fibres of $f$ are locally Noetherian.
The following are equivalent
\begin{enumerate}
\item $f$ is regular,
\item $f$ is flat and its fibres are geometrically regular schemes,
\item for every pair of affine opens $U \subset X$, $V \subset Y$
with $f(U) \subset V$ the ring map $\mathcal{O}_Y(V) \to \mathcal{O}_X(U)$
is regular,
\item there exists an open covering $Y = \bigcup_{j \in J} V_j$
and open coverings $f^{-1}(V_j) = \bigcup_{i \in I_j} U_i$ such
that each of the morphisms $U_i \to V_j$ is regular, and
\item there exists an affine open covering $Y = \bigcup_{j \in J} V_j$
and affine open coverings $f^{-1}(V_j) = \bigcup_{i \in I_j} U_i$ such
that the ring maps $\mathcal{O}_Y(V_j) \to \mathcal{O}_X(U_i)$ are regular.
\end{enumerate}
\end{lemma}

\begin{proof}
The equivalence of (1) and (2) is immediate from the definitions.
Let $x \in X$ with $y = f(x)$. By definition $f$ is flat at $x$
if and only if $\mathcal{O}_{Y, y} \to \mathcal{O}_{X, x}$ is a
flat ring map, and $X_y$ is geometrically regular at $x$ over
$\kappa(y)$ if and only if
$\mathcal{O}_{X_y, x} = \mathcal{O}_{X, x}/\mathfrak m_y\mathcal{O}_{X, x}$
is a geometrically regular algebra over $\kappa(y)$. Hence
Whether or not $f$ is regular at $x$
depends only on the local homomorphism of local rings
$\mathcal{O}_{Y, y} \to \mathcal{O}_{X, x}$.
Thus the equivalence of (1) and (4) is clear.

\medskip\noindent
Recall (More on Algebra, Definition \ref{more-algebra-definition-regular})
that a ring map $A \to B$ is regular if and only if it is flat
and the fibre rings $B \otimes_A \kappa(\mathfrak p)$ are Noetherian
and geometrically regular for all primes $\mathfrak p \subset A$.
By Varieties, Lemma \ref{varieties-lemma-geometrically-regular}
this is equivalent to $\Spec(B \otimes_A \kappa(\mathfrak p))$
being a geometrically regular scheme over $\kappa(\mathfrak p)$.
Thus we see that (2) implies (3). It is clear that (3) implies
(5). Finally, assume (5). This implies that $f$ is flat
(see Morphisms, Lemma \ref{morphisms-lemma-flat-characterize}).
Moreover, if $y \in Y$, then $y \in V_j$ for some $j$ and we see
that $X_y = \bigcup_{i \in I_j} U_{i, y}$ with each $U_{i, y}$
geometrically regular over $\kappa(y)$ by
Varieties, Lemma \ref{varieties-lemma-geometrically-regular}.
Another application of
Varieties, Lemma \ref{varieties-lemma-geometrically-regular}
shows that $X_y$ is geometrically regular. Hence (2) holds
and the proof of the lemma is finished.
\end{proof}

\end{document}
